% ============================================================
% 分片：s09-other  对应原文 PDF 第 111--119 页（印刷页 111--119）
% 负责章节：§9 Other groups + 9.1--9.4（§9 标题起）
% 译者：trans-s07   审校状态：初稿
% ============================================================

% ==== 原文 PDF 第 111 页（印刷页 111）====
% [跨页接续: 原文第 111 页开头两行仍属 §8（由 s08 分片负责），本分片自该页 §9 标题起]

\section{其他群}

\begin{keytakeaway}
\textbf{核心物理图像：}当 ${}^{L}G$-局部系统 $\mathcal{E}$ 有非平凡有限自同构群 $\Gamma$ 时，它在 $\mathcal{M}_H({}^{L}G)$ 中是轨形点；镜像对偶的 A-膜支撑在对偶奇异 Hitchin 纤维 $\boldsymbol{F}_b$ 上，并按 $\Gamma$ 的不可约表示分裂为分数本征膜——这就是一般群的内窥。
\textbf{数学工具：}Hitchin 纤维化与同调镜像对称、支撑在轨形点的凝聚层范畴 $\mathrm{Rep}(\Gamma)$、Wilson 与 't Hooft 算子、Ngô 的广义 Prym 簇 $P_b$ 与轨迹 $\boldsymbol{B}^{\mathrm{ani}}$、有限群上的 Fourier 变换、Grothendieck 的 faisceaux--fonctions 词典。
\textbf{主要结论：}正则本征膜/本征层是分数本征膜/本征层的直和；猜想 1 断言当 $b\in\boldsymbol{B}^{\mathrm{ani}}$ 时 $\Gamma$ 阿贝尔且 $\widehat{\Gamma}$ 是 $\pi_0(P_b)$ 的商；在函数域情形，用 $\Gamma$ 上的 Fourier 变换可从通常 Hecke 本征函数构造出满足分数 Hecke 性质的函数，从而验证上述猜想。
\end{keytakeaway}

本节把我们此前的结果与猜想概述性地推广到任意半单、单连通的 Lie 群 $G$ 的情形（后一假设并非本质的，只是为了简化叙述）。此时 ${}^{L}G$ 是伴随型的半单 Lie 群。

\subsection{概述}

回忆我们有两个对偶的 Higgs 丛模空间 $\mathcal{M}_H(G)$ 与 $\mathcal{M}_H({}^{L}G)$，以及相应的对偶 Hitchin 纤维化 (2.2)。几何 Langlands 对应在 [KW] 中被解释为这两个模空间之间的同调镜像对称，后者在一般纤维（光滑对偶环面）上归结为逐纤维的 T-对偶。在这一镜像对称下，$\mathcal{M}_H(G)$ 上 B-膜范畴与 $\mathcal{M}_H({}^{L}G)$ 上 A-膜范畴应当等价。我们关心的是 $\mathcal{M}_H({}^{L}G)$ 上对应于支撑在其轨形奇点的 B-膜的 A-膜。这样的奇点可以看作一个 ${}^{L}G$-局部系统 $\mathcal{E}$，其自同构群 $\Gamma$ 非平凡但有限。我们称这样的局部系统为「椭圆内窥的」，或为简短起见「内窥的」。于是 $\mathcal{E}$ 约化到一个或多个对偶（椭圆）内窥子群 ${}^{L}H\subset{}^{L}G$，后者定义为 $\Gamma$ 的非平凡元素的中心化子。

支撑在这样的点 $\mathcal{E}$ 上的 B-膜（等价地，凝聚层）范畴等价于 $\Gamma$ 的表示范畴 $\mathrm{Rep}(\Gamma)$。相应的 A-膜范畴的对象支撑在 $\mathcal{M}_H(G)$ 的 Hitchin 纤维 $\boldsymbol{F}_b$ 上，该纤维对偶于 $\mathcal{M}_H({}^{L}G)$ 中包含点 $\mathcal{E}$ 的 Hitchin 纤维 ${}^{L}\boldsymbol{F}_b$。因此 $b$ 是 $\mathcal{E}\in\mathcal{M}_H({}^{L}G)$ 在 Hitchin 底 $\boldsymbol{B}$ 中的像。本节要处理的问题是：描述对应于椭圆内窥 ${}^{L}G$-局部系统的 A-膜范畴，以及它们在 't Hooft/Hecke 算子作用下的性质。

在前几节中我们详细分析了 $G=SL_2$ 的情形。此时唯一对偶的椭圆内窥子群是 $O_2\subset SO_3={}^{L}G$，而约化到 $O_2$ 的一般局部系统具有自同构群 $\Gamma=\mathbb{Z}_2$。相应的 B-膜范畴等价于 $\mathrm{Rep}(\mathbb{Z}_2)$。\footnote{由于下文 §10 所述 gerbe 的扭转，这个等价并非典范的，但眼下我们忽略这一问题。}在 A-模型一侧，这对应于 Hitchin 纤维 $\boldsymbol{F}_b$ 有两个不可约分支这一事实。因此对偶的 A-膜范畴有两个支撑在这些分支上的不可约对象。这些分数 A-膜有附加参数；即秩一酉局部系统，它们对应于 $O_2$-局部系统。于是，我们得到对应于同态 $O_2\to SO_3$ 的转移（亦称函子性原理，或在物理解释下称为畴壁现象）在几何设定中的具体实现，如 §5.2.6 所解释。最后，这两个分数 A-膜

% ==== 原文 PDF 第 112 页（印刷页 112）====

满足分数 Hecke 性质，如 §5.3 所解释。在 §8 中，当曲线 $C$ 定义在有限域上时，我们把关于 A-膜的这一结果用 $\Bun_{SL_2}$ 上相应的 D-模与自守函数来解释（见 (1.1)）。本节我们提出这一图景的推广。

\subsection{对应于内窥局部系统的膜范畴}

设 $G$ 是半单单连通的复 Lie 群，${}^{L}G$ 是其 Langlands 对偶群（伴随型，中心平凡）。

若 ${}^{L}G$-局部系统 $\mathcal{E}$ 的自同构群是非平凡有限群，则称之为椭圆内窥的，或简称内窥的；我们把这个群记作 $\Gamma$。与此局部系统相联系的对偶内窥群 ${}^{L}H_s$\footnote{这里我们遵循传统，把内窥群称作 $H_s$ 而不是 ${}^{L}H_s$。}按定义是 $\Gamma$ 的非平凡元素 $s\in\Gamma$（$s\neq 1$）的中心化子。$\mathcal{E}$ 的结构群可以约化到任一 ${}^{L}H_s$。正如我们在 §5.2.6 已指出的，对任意子群 ${}^{L}H\subset{}^{L}G$ 都有自然的嵌入 $\mathcal{M}_H({}^{L}H)\subset\mathcal{M}_H({}^{L}G)$。因此我们看到，内窥局部系统 $\mathcal{E}$ 位于诸像 $\mathcal{M}_H({}^{L}H_s)$（$s\in\Gamma$，$s\neq 1$）在 $\mathcal{M}_H({}^{L}G)$ 中的像的交集里。

注意 $\mathcal{E}\in\mathcal{M}_H({}^{L}G)$ 也可以看作（在复结构 $I$ 下的）一个 Higgs 丛。它作为 Higgs 丛的自同构群同样是 $\Gamma$，且这个 Higgs 丛约化到子群 ${}^{L}H_s$，$s\in\Gamma$。

作为模空间 $\mathcal{M}_H({}^{L}G)$ 的点，内窥局部系统 $\mathcal{E}$ 是一个轨形点。记 $\mathcal{B}\text{-branes}_{\mathcal{E}}$ 为支撑在 $\mathcal{E}$ 的 B-膜（凝聚层）范畴。这个范畴等价于 $\mathrm{Rep}(\Gamma)$，不过等价可能并非典范的，如 §10 所述（在此情形，我们选取这样一个等价）。于是对 $\Gamma$ 的每个表示 $R$，我们在范畴 $\mathcal{B}\text{-branes}_{\mathcal{E}}$ 中有一个 B-膜 $\mathcal{B}_R$。按与 §5.3.2 相同的方式，我们发现 Wilson 算子 $W_{V,p}$（其中 $V\in\mathrm{Rep}({}^{L}G)$，$p\in C$）在这些膜上的作用由公式
\begin{equation}\tag{9.1}
W_{V,p}\cdot\mathcal{B}_R = \sum_{R'\in\mathrm{Irrep}(\Gamma)}V(R')_{\mathcal{E}_p}\otimes\mathcal{B}_{R'\otimes R},
\qquad R\in\mathrm{Rep}(\Gamma)
\end{equation}
给出。这里我们用了 $V$ 关于 ${}^{L}G\times\Gamma$ 的作用的分解
\begin{equation}\tag{9.2}
V = \bigoplus_{R'\in\mathrm{Irrep}(\Gamma)}V(R')\otimes R',
\end{equation}
其中 $\mathrm{Irrep}(\Gamma)$ 是 $\Gamma$ 的不可约表示的等价类之集。此外，对 ${}^{L}G$ 的任意表示 $U$，我们使用记号
\[
U_{\mathcal{E}_p} = \mathcal{E}_p\times_{{}^{L}G}U,
\]
其中 $\mathcal{E}_p$ 是 ${}^{L}G$-主齐性空间（torsor），即 $\mathcal{E}$ 在 $p\in C$ 处的纤维。

公式 (9.1) 表明，Wilson 算子的本征膜是正则表示
\[
\mathrm{Reg}(\Gamma) = \bigoplus_{R\in\mathrm{Irrep}({}^{L}G)}R^*\otimes R
\]
% [原文疑误: 此处以及 (9.3)、(9.5)、(9.8) 中求和指标均印作 $\mathrm{Irrep}({}^{L}G)$，但被求和的 $R$ 必须是 $\Gamma$ 的不可约表示，否则 $R^*\otimes R$、$\mathcal{B}_R$、$\mathcal{A}_R$、$\mathcal{F}_R$ 都没有意义；应为 $\mathrm{Irrep}(\Gamma)$。译文照录原文。 来源：原书 p.112 英文原文与译本对照（A）。]

% [Tutor Note: 推导补足。(9.1) 的实质是 $\Gamma$ 的表示环与 ${}^{L}G$ 的表示环的联合作用：Wilson 算子 $W_{V,p}$ 只通过 $V$ 作为 ${}^{L}G\times\Gamma$-表示的信息作用于 B-膜，分解 (9.2) 把 $V$ 拆成 $\bigoplus_{R'}V(R')\otimes R'$，于是作用在 $\mathcal{B}_R$ 上就得到 $V(R')_{\mathcal{E}_p}\otimes\mathcal{B}_{R'\otimes R}$。因此 $W_{V,p}$ 的本征膜必须对所有 $R'$ 同时对角化，而张量积 $R'\otimes R$ 正是 $\Gamma$ 表示环中的乘法。这也解释了为什么正则表示出现：只有它把每个不可约表示以重数 $\dim R$ 均匀地包含进来。 来源：原书 p.112 上下文（A）。]

% ==== 原文 PDF 第 113 页（印刷页 113）====

所对应的 B-膜的若干拷贝的直和，其中 $R^*$ 是 $R$ 的对偶。这个 B-膜是
\begin{equation}\tag{9.3}
\mathcal{B}_{\mathrm{Reg}(\Gamma)} = \bigoplus_{R\in\mathrm{Irrep}({}^{L}G)}R^*\otimes\mathcal{B}_R.
\end{equation}
我们有
\begin{equation}\tag{9.4}
W_{V,p}\cdot\mathcal{B}_{\mathrm{Reg}(\Gamma)} = V_{\mathcal{E}_p}\otimes\mathcal{B}_{\mathrm{Reg}(\Gamma)}.
\end{equation}
% [Tutor Note: 物理图像。(9.4) 说明 $\mathcal{B}_{\mathrm{Reg}(\Gamma)}$ 是 Wilson 算子的本征膜，其本征值是 $V_{\mathcal{E}_p}$，即局部系统 $\mathcal{E}$ 在 $p$ 处纤维上的向量空间；注意这里没有算子作用在 $V$ 上，因为本征值就是「把 $V$ 的纤维张量进来」。镜像对称把 Wilson 本征性翻译成 't Hooft 本征性，故下一段的 $\mathcal{A}_{\mathrm{Reg}(\Gamma)}$ 满足同形的 (9.6)，只是算子换成 $T_{V,p}$。 来源：原书 p.113 上下文（A）。]

现在考虑 $\mathcal{M}_H(G)$ 上 A-膜的镜像对偶范畴 $\mathcal{A}\text{-branes}_{\mathcal{E}}$。与前面几节中详细解释的 $G=SL_2$ 情形类比，我们预期这些 A-膜支撑在 $\mathcal{M}_H(G)$ 的 Hitchin 纤维 $\boldsymbol{F}_b$ 上，其中 $b$ 是 $\mathcal{E}\in\mathcal{M}_H({}^{L}G)$ 在 Hitchin 底 $\boldsymbol{B}$ 中的像。在 $G=SL_2$ 的情形我们看到，对内窥局部系统 $\mathcal{E}$，不同的 A-膜对应于 $\boldsymbol{F}_b$ 的不同分支。我们想理解一般情形会发生什么。

一个复杂因素是：Hitchin 纤维 ${}^{L}\boldsymbol{F}_b$ 上很可能存在对应于自同构群无限的 ${}^{L}G$-局部系统的点。一个例子是 $0\in\boldsymbol{B}$ 处的零纤维 ${}^{L}\boldsymbol{F}_0$，即幂零锥。这个纤维的点可以有有限或无限的自同构群。例如，对 ${}^{L}G=SO_3$，它包含平凡局部系统（此时 $\Gamma=SO_3$），也包含约化到子群 $\mathbb{Z}_2\times\mathbb{Z}_2$ 的局部系统（此时 $\Gamma=\mathbb{Z}_2\times\mathbb{Z}_2$），以及不可约局部系统。

若 Hitchin 纤维 ${}^{L}\boldsymbol{F}_b$ 包含自同构群无限的局部系统，则我们不能期望对偶 Hitchin 纤维 $\boldsymbol{F}_b$ 的结构由 ${}^{L}\boldsymbol{F}_b$ 的内窥点控制。事实上，同一个 $\boldsymbol{F}_b$ 会承载与支撑在这些局部系统上的 B-膜范畴镜像对偶的 A-膜范畴的对象。因此在这种情况下 $\boldsymbol{F}_b$ 的结构应当更为复杂。我们希望另文讨论这一更一般的情形，但眼下我们只限于不出现无限自同构群的情形。

对 $SL_n$，确保自同构群有限的一个有用条件是：谱曲线既约且不可约。这等价于要求 Higgs 场 $\varphi$ 的特征多项式不可约（这特别蕴含对一般的 $x\in C$，$\varphi(x)$ 是正则半单的）。这一判据对任意 $G$ 有简单的类比，\footnote{把 $\varphi$ 看作伴随表示中的矩阵，并令 $P(y)=\det(y-\varphi)$。对秩为 $r$ 的单群 $G$，一般地有 $P(y)=y^rQ(r)$（若 $Q$ 单链）；我们要的条件就是 $Q$ 不可约。对非单链的 $G$，一般地有 $P(y)=y^rQ(r)R(r)$，其中 $Q(r)$ 与 $R(r)$ 分别是长根与短根的贡献。在此情形，条件是 $Q$（或等价地 $R$）应当不可约。注意若 $(E,\varphi)$ 是使 $\varphi$ 满足这一不可约性判据的 Higgs 丛，则 $(E,\varphi)$ 自动稳定；$E$ 没有非平凡的 $\varphi$-不变子层，因而没有破坏稳定性的子层。于是在这一轨迹上，可以在只处理稳定 Higgs 丛的情况下应用 Hecke 算子。}
% [原文疑误: 脚注中 $P(y)=y^rQ(r)$ 以及 $Q(r)$、$R(r)$ 应为 $Q(y)$、$R(y)$：$P$ 是变量 $y$ 的多项式，而 $r$ 是 $G$ 的秩（常数），$Q(r)$ 与特征多项式的分解无关。译文照录原文。 来源：原书 p.113 英文原文与译本对照（A）。]
但它比确保自同构群有限所需要的条件更强。较弱的判据由 Ngô 在 [N1] 定义 7.5 中给出。由 [N1] 推论 7.6，它等价于下述条件。设 $P_b$ 是与 $b$ 相联系的广义 Prym 簇，定义于 [N1] 第 4 节（当 $G=SL_2$ 时，这就是 §5.2.2 所讨论的与 $b$ 相联系的谱曲线的 Prym 簇）。则条件是：$P_b$ 的分量群 $\pi_0(P_b)$ 有限。遵循 [N2]，我们把 $\boldsymbol{B}$ 中相应的轨迹记作 $\boldsymbol{B}^{\mathrm{ani}}$。

% ==== 原文 PDF 第 114 页（印刷页 114）====

设 $b\in\boldsymbol{B}^{\mathrm{ani}}$。此时 Hitchin 纤维是既约的，并包含一个开稠子集，它是阿贝尔群 $P_b$ 上的一个 torsor [N1]。于是 $\boldsymbol{F}_b$ 的不可约分支之集是群 $\pi_0(P_b)$ 上的一个 torsor，且 $\boldsymbol{F}_b$ 的所有分支重数为一。群 $\pi_0(P_b)$ 有限（由 $b\in\boldsymbol{B}^{\mathrm{ani}}$ 的假设）且阿贝尔。

现在从 [KW] 回忆：对一般的 $b\in\boldsymbol{B}$，对应于 $\boldsymbol{F}_b$（它是光滑环面）上任一秩一酉局部系统的 A-膜是对偶于作用在 B-膜上的 Wilson 算子的 't Hooft 算子的本征膜（见 §5.3）。我们猜想对任意 $b\in\boldsymbol{B}^{\mathrm{ani}}$ 这也成立。此外，我们猜想：在等价 $\mathcal{A}\text{-branes}_{\mathcal{E}}\simeq\mathcal{B}\text{-branes}_{\mathcal{E}}$ 下，对应于 $\mathcal{B}\text{-branes}_{\mathcal{E}}$ 的不可约对象 $\mathcal{B}_R$ 的 $\mathcal{A}\text{-branes}_{\mathcal{E}}$ 的每个不可约对象 $\mathcal{A}_R$ 都支撑在 $\boldsymbol{F}_b$ 的不可约分支的并上。

特别地，设 $\Gamma=\Gamma_b$ 是对偶 Hitchin 纤维 ${}^{L}\boldsymbol{F}_b$ 中局部系统可能的最大自同构群。那么自然预期：每个 $\mathcal{A}_R$ 支撑在 $\boldsymbol{F}_b$ 的某个特定不可约分支上，并且 $\mathrm{Irrep}(\Gamma_b)$ 与 $\boldsymbol{F}_b$ 的不可约分支之集（从而与集合 $\pi_0(P_b)$）之间存在一个双射（或许是既非典范的，如 $G=SL_2$ 时）。但 $\pi_0(P_b)$ 是阿贝尔群。这提示 $\mathrm{Irrep}(\Gamma_b)$ 也具有自然的阿贝尔群结构，且 $\Gamma_b$ 事实上是对偶于 $\pi_0(P_b)$ 的阿贝尔群。这样我们得到下述猜想。

\begin{namedconj}{1}
设 $\mathcal{E}$ 是椭圆内窥 ${}^{L}G$-局部系统，其自同构群为 $\Gamma$，并使 $\mathcal{E}$ 在 $\boldsymbol{B}$ 中的像 $b$ 属于 $\boldsymbol{B}^{\mathrm{ani}}$。则群 $\Gamma$ 是阿贝尔的，且其对偶群 $\widehat{\Gamma}$ 可以等同于与 $b$ 相联系的广义 Prym 簇 $P_b$ 的分量群 $\pi_0(P_b)$ 的一个商。

此外，若 $\Gamma_b$ 是对偶 Hitchin 纤维 ${}^{L}\boldsymbol{F}_b$ 中局部系统的最大自同构群，则 $\Gamma_b$ 同构于 $\pi_0(P_b)$ 的对偶群。
\end{namedconj}

% [Tutor Note: 推导补足。为什么 $b\in\boldsymbol{B}^{\mathrm{ani}}$ 会迫使 $\Gamma$ 阿贝尔？Ngô 的定理给出：在 $\boldsymbol{B}^{\mathrm{ani}}$ 上 Hitchin 纤维既约、所有不可约分支重数为一，且分支之集是 $\pi_0(P_b)$ 的 torsor（因为 $\boldsymbol{F}_b$ 的开稠子集是 $P_b$ 的 torsor）。若 $\mathrm{Irrep}(\Gamma_b)$ 要与分支之集（从而与 $\pi_0(P_b)$）双射，则 $\mathrm{Irrep}(\Gamma_b)$ 继承 $\pi_0(P_b)$ 的阿贝尔群结构；有限群的不可约表示集带有与群结构相容的阿贝尔群结构，当且仅当该群阿贝尔，故 $\Gamma_b$ 阿贝尔且 $\Gamma_b\cong\widehat{\pi_0(P_b)}$。这正是猜想 1 的推理，也是 Ngô 基本引理证明中 $\pi_0(P_b)$ 的 $\kappa$-分解与内窥群对应的几何版本。 来源：原书 p.114 上下文（A）。]

§5 中给出的结果在 $G=SL_2$ 的情形证实了这一猜想（另见第 10 页的脚注）。

\subsection{分数本征膜与本征层}

设 $\mathcal{E}$ 是内窥 ${}^{L}G$-局部系统。我们有上一小节讨论的镜像对偶范畴 $\mathcal{B}\text{-branes}_{\mathcal{E}}$ 与 $\mathcal{A}\text{-branes}_{\mathcal{E}}$。前者等价于 $\mathrm{Rep}(\Gamma)$，并包含与 $\Gamma$ 的不可约表示 $R$ 相联系的不可约对象 $\mathcal{B}_R$。相应的 A-膜记作 $\mathcal{A}_R$。因此 B-膜 (9.3) 对应于 A-膜
\begin{equation}\tag{9.5}
\mathcal{A}_{\mathrm{Reg}(\Gamma)} = \bigoplus_{R\in\mathrm{Irrep}({}^{L}G)}R^*\otimes\mathcal{A}_R.
\end{equation}
鉴于猜想 1，当 $\mathcal{E}$ 到 $\boldsymbol{B}$ 的投影属于 $\boldsymbol{B}^{\mathrm{ani}}$ 时，这一分解应当反映 Hitchin 纤维 $\boldsymbol{F}_b$ 分解为（若干）不可约分支之并。

由于 B-膜 $\mathcal{B}_{\mathrm{Reg}(\Gamma)}$ 是 Wilson 算子的本征膜（见公式 (9.4)），$\mathcal{A}_{\mathrm{Reg}(\Gamma)}$ 应当是 't Hooft 算子的本征膜：
\begin{equation}\tag{9.6}
T_{V,p}\cdot\mathcal{A}_{\mathrm{Reg}(\Gamma)} = V_{\mathcal{E}_p}\otimes\mathcal{A}_{\mathrm{Reg}(\Gamma)}.
\end{equation}

% ==== 原文 PDF 第 115 页（印刷页 115）====

我们称 $\mathcal{A}_{\mathrm{Reg}(\Gamma)}$ 为正则本征膜。

此外，公式 (9.1) 关于 Wilson 算子作用在 B-膜 $\mathcal{B}_R$（$R\in\mathrm{Irrep}(\Gamma)$）上的陈述，蕴含 't Hooft 算子 $T_{V,p}$（$V\in\mathrm{Rep}({}^{L}G)$）作用在相应 A-膜 $\mathcal{A}_R$ 上的下述公式：
\begin{equation}\tag{9.7}
T_{V,p}\cdot\mathcal{A}_R = \sum_{R'\in\mathrm{Irrep}(\Gamma)}V(R')_{\mathcal{E}_p}\otimes\mathcal{A}_{R'\otimes R},
\qquad R\in\mathrm{Rep}(\Gamma).
\end{equation}
我们称 A-膜 $\mathcal{A}_R$ 为分数本征膜。

注意我们预期公式 (9.6) 与 (9.7) 成立，无论 $\mathcal{E}$ 到 $\boldsymbol{B}$ 的投影是否属于 $\boldsymbol{B}^{\mathrm{ani}}$。若属于，则我们预期 $\Gamma$ 是阿贝尔的（见猜想 1）；否则它完全可能是非阿贝尔的，这可以从具体例子看出。

正如 §2.3 所解释的，我们预期 $\mathcal{M}_H(G)$ 上每个 A-膜都可以联系 $\Bun_G$ 上的一个 D-模。此外，A-膜在 't Hooft 算子作用下的性质，应当翻译为相应 D-模在 Hecke 算子作用下的类似性质。因此我们预言：$\Bun_G$ 上以 $\mathcal{E}$（一个内窥局部系统）为本征值的任一 Hecke 本征层，都是下述 D-模的若干拷贝的直和：
\begin{equation}\tag{9.8}
\mathcal{F}_{\mathrm{Reg}(\Gamma)} = \bigoplus_{R\in\mathrm{Irrep}({}^{L}G)}R^*\otimes\mathcal{F}_R,
\end{equation}
它满足正则 Hecke 性质
\[
H_V(\mathcal{F}_{\mathrm{Reg}(\Gamma)})\simeq V_{\mathcal{E}}\boxtimes\mathcal{F}_{\mathrm{Reg}(\Gamma)}
\]
（见公式 (6.5)）。此外，我们预言其构成项 $\mathcal{F}_R$ 是 $\Bun_G$ 上的不可约 D-模，它们满足公式 (9.7) 的类比
\begin{equation}\tag{9.9}
H_V(\mathcal{F}_R)\simeq\sum_{R'\in\mathrm{Irrep}(\Gamma)}V(R')_{\mathcal{E}}\boxtimes\mathcal{F}_{R'\otimes R},
\qquad R\in\mathrm{Rep}(\Gamma).
\end{equation}
这是公式 (6.10) 的一个变体，它意味着 D-模 $\mathcal{F}_R$（$R\in\mathrm{Irrep}(\Gamma)$）满足分数 Hecke 性质，从而是分数 Hecke 本征层。当 $G=SL_2$ 时，这些就是 §6.5 所讨论的 D-模 $\mathcal{F}_\pm$，且公式 (9.9) 与公式 (6.11) 一致。

于是，我们得到关于对应于内窥局部系统的（正则）Hecke 本征层结构的明确猜想：它们分裂为满足分数 Hecke 性质 (9.9) 的不可约 D-模的直和。我们是从对偶 Hitchin 纤维化的同调镜像对称、利用从 A-膜到 D-模的过渡导出这一猜想的。另一种看法是：从非阿贝尔版本的 Fourier--Mukai 变换 [La4, Ro] 的角度来看，这一变换由 A. Beilinson 与 V. Drinfeld 提出，应当是 $C$ 上 ${}^{L}G$-局部系统模叠上的 $\mathcal{O}$-模范畴与 $\Bun_G$ 上 D-模范畴之间某个范畴等价（其精确定义目前尚不清楚）（例如见 [F1] 第 6.2 节）。

% ==== 原文 PDF 第 116 页（印刷页 116）====

\subsection{Hecke 本征函数的计算}

上一节中我们对对应于内窥局部系统的 Hecke 本征层的结构作了猜想。迄今为止，我们讨论的是复曲线 $C$ 上的局部系统。然而我们猜想：若改为考虑定义在有限域上的曲线上的 $\ell$-进局部系统，或等价地，考虑 $\ell$-进同态 $\sigma: W_F\to{}^{L}G$（其中 $W_F$ 是该曲线的函数域 $F$ 的 Weil 群），同样的模式仍然成立。此时群 $\Gamma$ 的类比是 $\sigma$ 的像的中心化子，传统上记作 $S_\sigma$。但这里我们沿用同一记号 $\Gamma$。在此语境下有一个新要素，即 Grothendieck 的 faisceaux--fonctions 词典。它使我们能从 Hecke 本征层（现在视为 $\Bun_G$ 上的反常层）过渡到 adele 群 $G(\mathbb{A}_F)$ 的某个双商上相应的自守函数，从而有机会检验我们的猜想。我们已经对 $G=SL_2$ 在 §8.6 中做了这件事，并表明这样的函数确实存在。这里我们把分析推广到上面所考虑的一般情形。

\subsubsection{阿贝尔情形}

设我们有一个内窥同态 $\sigma: W_F\to{}^{L}G$。这意味着其像的中心化子 $\Gamma$ 是非平凡的有限群（我们仍假定 ${}^{L}G$ 是伴随型半单群）。让我们从 $\Gamma$ 阿贝尔的情形开始。回忆在 $\mathbb{C}$ 上的类似情形中，我们预期不可约 D-模由 $\Gamma$ 的一维表示（特征标）标记（或许略非典范地）。于是，对每个 $\chi\in\widehat{\Gamma}=\mathrm{Irrep}(\Gamma)$，我们有一个 D-模 $\mathcal{F}_\chi$。它们必须满足分数 Hecke 性质 (9.9)。过渡到有限域上的曲线后，我们应当有 $\Bun_G$ 上满足同一性质的相应反常层，并与自守函数 $f_\chi$ 相联系。层的分数 Hecke 性质翻译为这些函数上的下述方程：
\begin{equation}\tag{9.10}
T_{V,x}\cdot f_\chi = \sum_{\mu\in\widehat{\Gamma}}a_{V,\mu,x}f_{\chi\cdot\mu},
\qquad \chi\in\widehat{\Gamma}.
\end{equation}
这里 $V$ 是 ${}^{L}G$ 的表示，它在 $\Gamma$ 的作用下如下分解
\begin{equation}\tag{9.11}
V = \bigoplus_{\mu\in\widehat{\Gamma}}V(\mu),
\end{equation}
并且
\begin{equation}\tag{9.12}
a_{V,\mu,x} = \Tr(\sigma(\mathrm{Fr}_x),V(\mu)),
\end{equation}
其中 $\sigma: W_F\to{}^{L}G$ 是代替局部系统 $\mathcal{E}$ 的对象（由于按我们的假设 $\sigma$ 落在 $\Gamma$ 的中心化子中，右端是良定义的）。$T_{V,x}$ 是对应于闭点 $x\in C$ 的经典 Hecke 算子（见 §7.1）。

现在我们将证明满足分数 Hecke 性质 (9.10) 的函数 $f_\chi$ 确实存在，并且可以通过在 $\Gamma$ 上作某种 Fourier 变换从通常的 Hecke 本征函数得到。这将推广 §8.6 在 $G=SL_2$ 情形所得的公式。

% ==== 原文 PDF 第 117 页（印刷页 117）====

为简单起见，这里假定 $x$ 是 $C$ 的剩余域等于基域的闭点。对其他闭点，计算需按 §8.6 的方式修改。

由 (9.10)，在 $\Gamma$ 上作 Fourier 变换，我们得到 $T_{V,x}$ 的本征函数：
\begin{equation}\tag{9.13}
\widehat{f}_\gamma = \sum_{\chi\in\widehat{\Gamma}}\chi(\gamma)f_\chi,
\qquad \gamma\in\Gamma,
\end{equation}
其本征值为
\begin{equation}\tag{9.14}
A_{x,\gamma} = \sum_{\mu\in\widehat{\Gamma}}\mu(\gamma)a_{V,\mu,x}.
\end{equation}
特别地，
\[
A_{x,1} = \sum_{\mu\in\widehat{\Gamma}}a_{V,\mu,x} = \Tr(\sigma(\mathrm{Fr}_x),V),
\]
所以
\[
\widehat{f}_1 = \sum_{\chi\in\widehat{\Gamma}}f_\chi
\]
是相应于 $\sigma$ 的 Hecke 本征函数，正如所预期。

但其他函数 $\widehat{f}_\gamma$（$\gamma\neq 1$）呢？我们断言它们也是 Hecke 本征函数，但相应于其他同态
\[
\sigma_\gamma: W_F\to{}^{L}G.
\]
确切地说，回忆我们有同态
\[
\mathrm{res}: W_F\to W_{\mathbb{F}_q}=\mathbb{Z},
\]
它由限制到标量 $\mathbb{F}_q\subset F$ 得到。设
\[
\alpha_\gamma: W_F\to\Gamma
\]
是由 $\mathrm{res}$ 与同态 $\mathbb{Z}\to\Gamma$（$1\mapsto\gamma$）复合得到的同态。由于 $\Gamma$ 中心化 $\sigma$ 的像，公式
\[
\sigma_\gamma(g) = \sigma(g)\alpha_\gamma(g)
\]
对每个 $\gamma\in\Gamma$ 定义了同态 $W_F\to{}^{L}G$。我们断言
\[
A_{x,\gamma} = \Tr(\sigma_\gamma(\mathrm{Fr}_x),V),
\]
于是函数 $\widehat{f}_\gamma$ 事实上是相应于 $\sigma_\gamma: W_F\to{}^{L}G$ 的 Hecke 本征函数！（我们回忆，在上述计算中假定了 $x$ 是 $C$ 的 $\mathbb{F}_q$-点。对其他闭点，应用 §8.6 的分析可得同样结果。）

现在作逆 Fourier 变换，我们用通常的 Hecke 本征函数表达对应于层 $\mathcal{F}_\chi$ 的函数 $f_\chi$：
\begin{equation}\tag{9.15}
f_\chi = \frac{1}{|\Gamma|}\sum_{\gamma\in\Gamma}\chi(\gamma)\widehat{f}_\gamma.
\end{equation}
% [Tutor Note: 推导补足。Fourier 变换之所以奏效，是因为 (9.10) 中作用在函数 $f_\chi$ 上的核 $a_{V,\mu,x}$ 在 $\widehat{\Gamma}$ 上卷积；对 $\chi$ 作变换即把卷积对角化。(9.14) 表明本征值是 $\mu\mapsto a_{V,\mu,x}$ 在 $\gamma$ 处的求值，而它恰好等于 $\Tr(\sigma_\gamma(\mathrm{Fr}_x),V)$。关键是 $\sigma_\gamma = \sigma\cdot\alpha_\gamma$ 的构造：$\alpha_\gamma$ 通过 $\mathrm{res}: W_F\to W_{\mathbb{F}_q}=\mathbb{Z}$ 把「模 Dedekind 类」的 $\mathbb{Z}$ 分量与 $\Gamma$ 中的元素结合，因此它只改变各 Frobenius 的标量部分。于是每个 $\widehat{f}_\gamma$ 都是某个扭转同态 $\sigma_\gamma$ 的普通 Hecke 本征函数，而 $f_\chi$ 是它们按 $\Gamma$ 的 Fourier 变换得到的线性组合——这正是「几何上正确的对象不是本征函数而是本征函数的线性组合」的机制。 来源：原书 p.117 上下文（A）。]

% ==== 原文 PDF 第 118 页（印刷页 118）====

这推广了我们在 $SL_2$ 情形中的公式
\[
f_\pm = \frac{1}{2}(f_\sigma\pm f_{\sigma'})
\]
（见公式 (8.24)）。

满足分数 Hecke 性质 (9.10) 的函数论类比的函数 $f_\chi$ 的存在性（我们从 A-膜学到了这一点），为我们的预言提供了一致性检验。

总之：我们发现，几何上「正确」的对象（对应于不可约反常层 $\mathcal{F}_\chi$，即「分数」Hecke 本征层）不是通常的 Hecke 本征函数，而是它们的线性组合（由有限 Fourier 变换得到），后者对应于由 $\gamma\in\Gamma$ 标记的一族 Galois 表示 $\{\sigma_\gamma\}$。它们被构造为 $\sigma$ 的简单扭转。我们注意到公式 (9.15) 中的 Fourier 变换与 Lusztig 在特征层理论 [Lu] 中观察到的 Fourier 变换有些相似。

\subsubsection{非阿贝尔情形}

现在考虑 $\Gamma$ 非阿贝尔的情形。在 $\mathbb{C}$ 上，假定猜想 1，这意味着 Hitchin 模空间 $\mathcal{M}_H({}^{L}G)$ 中相应的点不是一般正则半单的。在此情形，我们预期对偶 Hitchin 纤维 $\boldsymbol{F}_b$ 的某些分支的重数大于 1，且应当等于 $\Gamma$ 的相应不可约表示 $R$ 的维数。

按 §9.3 的猜想，把它们搬到有限域上曲线的情形，我们有 $\Bun_G$ 上满足分数 Hecke 性质 (9.9) 的不可约反常层 $\mathcal{F}_R$（$R\in\mathrm{Irrep}(\Gamma)$）。设 $f_R$（$R\in\mathrm{Irrep}(\Gamma)$）是相应的自守函数。于是我们有公式 (9.10) 的类比
\begin{equation}\tag{9.16}
T_{V,x}\cdot f_R = \sum_{R'\in\mathrm{Irrep}(\Gamma)}a_{V,R',x}f_{R\otimes R'},
\qquad R\in\mathrm{Irrep}(\Gamma),
\end{equation}
其中我们使用分解 (9.2)，并令
\begin{equation}\tag{9.17}
a_{V,R',x} = \Tr(\sigma(\mathrm{Fr}_x),V(R')).
\end{equation}
如前，
\[
R\otimes R' = \bigoplus_{R''\in\mathrm{Irrep}(\Gamma)}(R'')^{\oplus m^{R,R'}_{R''}},
\]
且按定义
\[
f_{R\otimes R'} = \sum_{R''\in\mathrm{Irrep}(\Gamma)}m^{R,R'}_{R''}f_{R''}.
\]

让我们求 $T_{V,x}$ 的本征函数。它们由 $\Gamma$ 中的共轭类 $[\gamma]$ 标记，并由公式
\begin{equation}\tag{9.18}
\widehat{f}_{[\gamma]} = \sum_{R\in\mathrm{Irrep}(\Gamma)}\Tr([\gamma],R)f_R,
\qquad [\gamma]\in\Gamma.
\end{equation}
给出。相应的本征值是
\begin{equation}\tag{9.19}
A_{x,[\gamma]} = \sum_{R'\in\mathrm{Irrep}(\Gamma)}\Tr([\gamma],R')a_{V,R',x}.
\end{equation}
% [Tutor Note: 易错预警。非阿贝尔情形把阿贝尔情形的一维特征标 $\chi(\gamma)$ 换成中心特征标 $\Tr([\gamma],R)$：因为 $R$ 可能高维，$\Gamma$ 的不可约表示不再是特征标，只有共轭类上的迹函数才是良定义的。(9.18) 因此是对共轭类的 Fourier 变换，而不再是群元素上的 Fourier 变换；这也解释了为什么求和指标写作 $[\gamma]\in\Gamma$。$f_{R\otimes R'}$ 的多重数定义使 (9.16) 与 (9.10) 在形式上统一。 来源：原书 p.118 上下文（A）。]

% ==== 原文 PDF 第 119 页（印刷页 119）====

特别地，
\[
A_{x,[1]} = \sum_{R'\in\mathrm{Irrep}(\Gamma)}\dim(R')a_{V,R',x} = \Tr(\sigma(\mathrm{Fr}_x),V),
\]
所以
\[
\widehat{f}_{[1]} = \sum_{R\in\mathrm{Irrep}(\Gamma)}\dim(R)f_R
\]
是相应于 $\sigma$ 的 Hecke 本征函数，正如所预期。

其他函数 $\widehat{f}_{[\gamma]}$（$[\gamma]\neq[1]$）也是 Hecke 本征函数，但相应于其他同态
\[
\sigma_{[\gamma]}: W_F\to{}^{L}G,
\]
后者由公式
\[
\sigma_{[\gamma]}(g) = \sigma(g)\alpha_\gamma(g)
\]
定义，其中
\[
\alpha_\gamma: W_F\to\Gamma
\]
是由 $\mathrm{res}$（见公式 (8.26)）与同态 $\mathbb{Z}\to\Gamma$（$1\mapsto\gamma$）复合得到的同态，而 $\gamma$ 是 $[\gamma]$ 的任意元素。显然，$\sigma_{[\gamma]}$ 的等价类只依赖 $[\gamma]$，不依赖 $\gamma$ 的选取。

相应的 Hecke 本征值 (9.19) 是
\[
A_{x,[\gamma]} = \Tr(\sigma_\gamma(\mathrm{Fr}_x),V),
\]
于是函数 $\widehat{f}_\gamma$ 事实上是相应于 $\sigma_\gamma: W_F\to{}^{L}G$ 的 Hecke 本征函数，正如所愿。
% [存疑: 此处印作 $\widehat{f}_\gamma$，而上文非阿贝尔情形一律用 $\widehat{f}_{[\gamma]}$（$\gamma$ 只是共轭类的代表元）；按上下文应读作 $\widehat{f}_{[\gamma]}$。译文照录原文。 来源：原书 p.119 英文原文；存疑待核（A）。]

现在我们可以用通常的 Hecke 本征函数表达对应于层 $\mathcal{F}_R$ 的函数 $f_R$：
\[
f_\chi = \frac{1}{|\Gamma|}\sum_{[\gamma]\in\Gamma}\Tr([\gamma],R)\widehat{f}_{[\gamma]},
\]
如同阿贝尔情形。
% [原文疑误: 左端应为 $f_R$：本段讨论的是与层 $\mathcal{F}_R$ 相应的函数 $f_R$，且右端出现 $R$、上文记号亦为 $f_R$；此处印作 $f_\chi$ 系沿袭阿贝尔情形的记号。译文照录原文。 来源：原书 p.119 英文原文与译本对照（A）。]

于是，在非阿贝尔情形下满足分数 Hecke 性质 (9.16) 的函数也存在。这为 §9.3 的猜想提供了一致性检验。

% [以下为 §10 标题起的内容，属 s10-gerbes 分片]
